\chapter{Emphysema}
\index{Emphysema}

\section*{Definition and Overview}
Emphysema is a type of chronic obstructive pulmonary disease (COPD) in which the air sacs of the lungs, called alveoli, are progressively destroyed. This damage reduces the surface area available for oxygen exchange, and the lungs lose their elasticity, trapping air and making it difficult to breathe out fully. Emphysema develops gradually over many years, is almost always caused by smoking, and causes progressive breathlessness. There is no cure, but the disease can be slowed, symptoms managed, and quality of life maintained with good care.

\section*{Epidemiology}
Emphysema affects a substantial proportion of older adults and is one of the most common causes of chronic respiratory disability in many countries. It is strongly linked to smoking, and most people with emphysema are current or former smokers. The disease is more common in men, though the gap has narrowed as smoking patterns changed. Around one in four long-term smokers develops clinically significant COPD, of which emphysema is a major component.

\section*{Aetiology and Pathophysiology}
Emphysema results from long-term lung damage, most often from cigarette smoke.
\begin{itemize}
    \item \textbf{Smoking:} the overwhelming cause; cigarette smoke triggers chronic inflammation and activates enzymes that destroy the walls of the air sacs
    \item \textbf{Other exposures:} long-term exposure to dust, chemicals, and air pollution; in rare cases, an inherited deficiency of alpha-1 antitrypsin causes emphysema in non-smokers
    \item \textbf{Mechanism:} inflammation and protease enzymes break down the elastic walls of the alveoli
    \item \textbf{Effects:} the air sacs merge into large, inefficient spaces, reducing oxygen exchange; loss of elasticity traps air and flattens the diaphragm
    \item \textbf{Progression:} damage is irreversible and continues while the exposure continues
    \item \textbf{Hyperinflation:} trapped air over-inflates the lungs, making breathing effortful
\end{itemize}

\section*{Clinical Features}
Emphysema develops insidiously over years.
\begin{itemize}
    \item Gradually worsening breathlessness, first on exertion and later at rest
    \item A chronic cough, sometimes with phlegm
    \item Wheezing and chest tightness
    \item A "barrel chest" from hyperinflated lungs
    \item Prolonged breathing out, often through pursed lips
    \item Fatigue, weight loss, and loss of muscle mass in advanced disease
    \item Frequent chest infections
    \item Bluish lips and fingers with low oxygen in advanced disease
\end{itemize}

\section*{Differential Diagnosis}
Other conditions cause similar symptoms.
\begin{itemize}
    \item Chronic bronchitis, the other major form of COPD, with a prominent cough and phlegm
    \item Asthma, which is usually reversible and episodic
    \item Heart failure, which causes breathlessness and fluid retention
    \item Bronchiectasis, with chronic productive cough
    \item Lung cancer and other causes of breathlessness
\end{itemize}

\section*{When to Seek Medical Care}
Anyone with persistent breathlessness, a chronic cough, or a smoking history should be assessed by a GP. Once diagnosed, regular reviews are important to manage the disease and prevent flares.

\textbf{Contact your local emergency services for:}
\begin{itemize}
    \item Sudden severe breathlessness or breathlessness at rest
    \item Chest pain or signs of a heart attack
    \item Coughing up blood
    \item Bluish lips or confusion, which indicate very low oxygen
    \item Difficulty speaking full sentences
\end{itemize}

\section*{Diagnosis and Investigations}
Diagnosis is confirmed with breathing tests.
\begin{itemize}
    \item \textbf{Spirometry:} measures how much and how fast air is exhaled; a reduced ratio of forced expiratory volume to forced vital capacity confirms COPD
    \item \textbf{Chest X-ray:} may show hyperinflation and flattened diaphragm
    \item \textbf{CT scan:} shows the extent of the air sac damage
    \item \textbf{Oxygen measurement:} pulse oximetry and blood gas tests
    \item \textbf{Alpha-1 antitrypsin testing:} for non-smokers and younger people with emphysema
\end{itemize}

\section*{Management}
Management slows progression, relieves symptoms, and prevents flares.
\begin{itemize}
    \item \textbf{Smoking cessation:} the single most important step; stopping slows the decline in lung function
    \item \textbf{Inhalers:} bronchodilators open the airways, and inhaled corticosteroids reduce flare-ups in suitable people
    \item \textbf{Pulmonary rehabilitation:} exercise training and education that improve breathlessness and quality of life
    \item \textbf{Vaccinations:} influenza, pneumococcal, and COVID-19 vaccines prevent infections
    \item \textbf{Oxygen therapy:} home oxygen for people with low oxygen levels
    \item \textbf{Management of flare-ups:} antibiotics and corticosteroids for exacerbations
    \item \textbf{Surgery:} lung volume reduction and, in severe cases, lung transplantation
\end{itemize}

\section*{Complications}
\begin{itemize}
    \item Frequent exacerbations requiring hospital admission
    \item Respiratory failure and low oxygen
    \item Pulmonary hypertension and strain on the heart
    \item Osteoporosis from inactivity and steroid use
    \item Weight loss and muscle wasting
    \item Depression and anxiety
    \item Increased risk of lung cancer in smokers
\end{itemize}

\section*{Prognosis and Outcomes}
Emphysema is a progressive disease, but the rate of decline varies greatly and is strongly influenced by whether the person stops smoking. With smoking cessation, pulmonary rehabilitation, and optimal treatment, most people maintain a good quality of life for many years. The disease is incurable, but it is very treatable, and severe disability is not inevitable.

\section*{Prevention and Self-Care}
\begin{itemize}
    \item Do not smoke, and seek help to quit if you do
    \item Avoid exposure to dust, fumes, and second-hand smoke
    \item Take inhalers correctly and as prescribed
    \item Attend pulmonary rehabilitation
    \item Keep vaccinations up to date
    \item Maintain physical activity within your limits
    \item Eat well and maintain a healthy weight
\end{itemize}

\section*{Sources and Further Reading}
The following reputable sources provide further information for healthcare students and patients seeking reliable, current advice about emphysema.
\begin{itemize}
    \item Lung Foundation Australia. \textit{COPD and emphysema information.} https://lungfoundation.com.au/
    \item Healthdirect Australia. \textit{COPD and emphysema.} https://www.healthdirect.gov.au/copd
    \item NHS. \textit{Emphysema: overview and treatment.} https://www.nhs.uk/conditions/chronic-obstructive-pulmonary-disease-copd/
    \item Mayo Clinic. \textit{Emphysema: symptoms and causes.} https://www.mayoclinic.org/diseases-conditions/emphysema/
    \item MedlinePlus. \textit{Emphysema.} https://medlineplus.gov/emphysema.html
    \item MSD Manual. \textit{Emphysema: professional overview.} https://www.msdmanuals.com/
\end{itemize}
